% Visual Semantic Profiling of the Voynich Manuscript
% Jacob Lyons — xenoglyph — April 2026
%
% arXiv source bundle. Body content is \input{body.tex}, which is
% auto-generated from the canonical markdown source
% (../voynich_visual_semantics_preprint.md) by:
%   pandoc --natbib  →  postprocess.py  →  build/body.tex
% See Makefile for the pipeline. Everything above \begin{document} is
% hand-written and lives here so the preamble stays reviewable.

\documentclass[11pt]{article}

% ── Character encoding, fonts, language ───────────────────────────────
\usepackage[utf8]{inputenc}
\usepackage[T1]{fontenc}
% Map the unicode math characters the author used in prose (outside
% $..$ math mode) onto their LaTeX equivalents. Without these the
% source markdown's "≥ 84%" and "2770 × 3770" trip pdfTeX with
% "Unicode character ≥/× not set up for use with LaTeX". This covers
% every stray unicode math symbol found in the body — add more
% here if a future paper uses additional Unicode in prose.
\DeclareUnicodeCharacter{00D7}{\ensuremath{\times}}      % ×
\DeclareUnicodeCharacter{2265}{\ensuremath{\geq}}        % ≥
\DeclareUnicodeCharacter{2264}{\ensuremath{\leq}}        % ≤
\DeclareUnicodeCharacter{2212}{\ensuremath{-}}           % − (minus)
\DeclareUnicodeCharacter{2248}{\ensuremath{\approx}}     % ≈
\DeclareUnicodeCharacter{2260}{\ensuremath{\neq}}        % ≠
% lmodern provides a scalable outline family that microtype's font
% expansion needs — without it pdfTeX trips on "auto expansion is only
% possible with scalable fonts" because default Computer Modern is
% bitmap.
\usepackage{lmodern}
\usepackage[english]{babel}

% ── Math + science typography ─────────────────────────────────────────
\usepackage{amsmath}
\usepackage{amssymb}
\usepackage{mathtools}

% ── Figures ────────────────────────────────────────────────────────────
\usepackage{graphicx}
\graphicspath{{figures/}}
\usepackage{caption}
\captionsetup{font=small,labelfont=bf,labelsep=colon,singlelinecheck=false,
              justification=raggedright}

% ── Tables ─────────────────────────────────────────────────────────────
\usepackage{booktabs}
\usepackage{longtable}
\usepackage{array}
\usepackage{multirow}
% pandoc uses calc-style lengths like (\linewidth - 4\tabcolsep) * \real{.33}
% in longtable column specs — the \real macro needs this package.
\usepackage{calc}
\providecommand{\real}[1]{#1}

% ── Bibliography ──────────────────────────────────────────────────────
% natbib + plainnat gives the author-year style arXiv cs.CV reviewers
% expect, and matches the \citep / \citet pandoc already emits.
\usepackage[round,authoryear]{natbib}
\bibliographystyle{plainnat}

% ── Hyperlinks (loaded last, as hyperref requires) ────────────────────
\usepackage[unicode,breaklinks=true]{hyperref}
\hypersetup{
  colorlinks=true,
  linkcolor=[rgb]{0.0,0.2,0.5},
  citecolor=[rgb]{0.0,0.3,0.2},
  urlcolor=[rgb]{0.0,0.3,0.5},
  pdftitle={Visual Semantic Profiling of the Voynich Manuscript},
  pdfauthor={Jacob Lyons},
  pdfsubject={Computer Vision; Digital Humanities},
  pdfkeywords={Voynich, CLIP, visual semantics, zero-shot, cs.CV, cs.DL}
}
\usepackage{url}

% ── Page geometry ─────────────────────────────────────────────────────
\usepackage[margin=1.25in]{geometry}
\usepackage{setspace}

% ── Small niceties ────────────────────────────────────────────────────
\usepackage{xcolor}
\usepackage{microtype}        % better justification, fewer overfulls
\usepackage[inline]{enumitem} % pandoc's \tightlist macro (see below)
% pandoc emits \tightlist to compress markdown lists — define a no-op
% in case pandoc's header-includes isn't regenerated.
\providecommand{\tightlist}{\setlength{\itemsep}{0pt}\setlength{\parskip}{0pt}}

% pandoc also emits \passthrough{…} around inline-code ranges in some
% modes; no-op it so the text renders unchanged.
\providecommand{\passthrough}[1]{#1}

% pandoc suppresses longtable auto-captions by redefining \LTcaptype
% to "none" and then issuing \addtocounter{\LTcaptype}{-1}. LaTeX
% needs a matching counter named "none" in scope or the increment
% errors out. Zero cost, fully idiomatic pandoc workaround.
\newcounter{none}

% \real is a pandoc column-width helper that stock longtable doesn't
% know about. Define it as identity so the generated column specs
% compile.
\providecommand{\real}[1]{#1}

% ── Title block ───────────────────────────────────────────────────────
\title{Visual Semantic Profiling of the Voynich Manuscript:\\
       \large Reading Meaning from Illustrations in an Undeciphered Codex}
\author{Jacob Lyons\\
        \textit{xenoglyph}\\
        \href{https://github.com/xenoglyph-ai/voynich-public}{github.com/xenoglyph-ai/voynich-public}}
\date{April 2026}

\begin{document}
\maketitle

% ── Abstract ──────────────────────────────────────────────────────────
% Abstract mirrors the YAML frontmatter in the source markdown verbatim.
% arXiv's abstract field will be populated from this block.
\begin{abstract}
The Voynich Manuscript (Beinecke MS 408) is a 240-page illustrated codex
carbon-dated to 1404--1438 whose script has resisted decipherment for
more than six centuries. Every quantitative decipherment attempt on
record --- from William Friedman's WWII codebreakers to contemporary
neural language models --- has treated the manuscript as a text-first
object. We present the first systematic \emph{computational} visual
semantic analysis of the manuscript, treating its 206 page images as
the primary signal. Using a zero-shot visual semantic profiling platform
in which a frozen vision-language foundation model is scored against
sixteen natural-language archetype descriptors, we profile every
analysable page of the Beinecke digital facsimile (197 pages after
excluding nine covers and binding flyleaves) and measure whether the
resulting sixteen-dimensional profiles discriminate between the
scholarly section taxonomy. All sixteen dimensions discriminate between
sections at $p < 10^{-15}$ under one-way ANOVA, Welch's robust ANOVA,
and Kruskal--Wallis tests, with effect sizes ($\eta^2$) between
0.30 and 0.83. A multinomial logistic regression --- fit inside a
\texttt{Pipeline} that eliminates cross-validation scaler leakage ---
recovers scholarly section labels with 90.4\% leave-one-out accuracy
(Wilson 95\% CI [85.4\%, 93.7\%]; permutation-test $p < 10^{-3}$),
against a 20\% chance baseline and a 59.9\% majority-class baseline.
As ablations we report 92.4\% accuracy on the raw 768-d foundation-model
embeddings (the archetype projection trades a small amount of accuracy
for full interpretability) and 72.1\% on six handcrafted layout features
(layout alone carries real but strictly weaker signal). A
lens-specificity control shows that two alternative 16-d archetype
lenses --- a cross-cultural archaeological lens and a hermetic
cryptological lens --- \emph{also} recover section structure well above
chance (84.8\% and 87.3\% respectively), confirming that the recovered
signal is a property of the manuscript rather than of the
target-appropriate lens. A head-to-head comparison against a
character-n-gram classifier on Takeshi Takahashi's complete Voynichese
transcription shows that the text channel recovers sections at 92.3\%
--- statistically indistinguishable from the 16-d visual channel on
the same 182-page intersection --- while the raw 768-d visual embedding
achieves 96.2\%: both channels carry the section structure, and the
visual channel offers the unique advantage of producing profiles that
are human-readable in a known language. An out-of-distribution sanity
check applies the voynich lens via a local off-the-shelf CLIP pipeline
to the Tacuinum Sanitatis (a same-era medieval herbal peer), the Codex
Seraphinianus, and the Rohonc Codex: the lens fires on
herbal/pharmaceutical/fertility dimensions for the Tacuinum and on
\emph{encoded/hidden} for the two undeciphered-script codices,
confirming that the lens is a legitimate medieval-codex content
detector rather than a Voynich-specific artefact. The only class where
the Voynich classifier fails badly is the pharmaceutical section, which
is misclassified almost exclusively as \emph{herbal} with 50\% recall
but 91\% precision --- a failure that matches the scholarly observation
that pharmaceutical pages share the same plant specimens as the herbal
section, and we show the individual misclassified pages in a
qualitative gallery. We conclude that the illustrations of the Voynich
Manuscript encode structured thematic content accessible to
computational analysis even though the underlying text is not, that
the scholarly section structure is independently recoverable by a
non-textual method, and that the broader visual channel of the
manuscript is considerably richer than the text-first literature has
assumed. We do not claim to have read the manuscript. We claim only
that one more channel of the manuscript --- its visual channel --- is
not empty.
\end{abstract}

% ── Body (auto-generated from the markdown source) ────────────────────
\input{body}

% ── Bibliography ──────────────────────────────────────────────────────
\bibliography{references}

\end{document}
